\documentclass{article}
\usepackage[utf8]{inputenc}
\usepackage{graphicx}
\usepackage{multicol}
\usepackage{caption}
\usepackage{vmargin}
\usepackage[hidelinks]{hyperref}
\title{Laboratorio Perdidas de Energía Por Accesorios (Adimentados)}
\author{Lozada Rodriguez Jessica Lorena - Barrera Obando Camilo Enrique}
\date{28 de Abril de 2019}

\begin{document}
\begin{multicols}{2}
[
\maketitle
\section{Palabras Claves}
Perdidas primarias, Perdidas secundarias, Velocidad, Fricción, Longitud, Diametro
\section{Resumen}
La disminución de energía que se ocasiona  por accesorios en la tubería es llamada perdidas por accesorios pueden diferenciarse por codos, válvulas etc. En este informe de laboratorio tiene como propósito identificar, analizar y calcular las pérdidas por accesorios de un fluido en un sistema con tuberías y codos basándonos en la ecuación de energía de Bernoulli, los datos se tomaron en el laboratorio de Hidráulica de la Universidad la Gran Colombia
]
\section{Marco Teórico}
\textbf{TIPOS DE FLUJOS:}
Existen muchas maneras de clasificar el movimiento de los fluidos, de acuerdo con la estructura del flujo, o con la situación o configuración física. Ejemplo: 
flujo externo, flujo interno, flujo permanente, flujo uniforme, flujo compresible sin rozamiento.

\begin{itemize}
\item \textbf{Conservación de la energía. Ecuación de Bernoulli:}
\end{itemize}
El teorema de Bernoulli implica una relación entre los efectos de la presión, la velocidad y la gravedad, e indica que la velocidad aumenta cuando la presión disminuye. 
\includegraphics[scale=0.45]{bernoulli.JPG}
\begin{itemize}
\item \textbf{pérdida de carga por accesorios }
\end{itemize}
Las pérdidas de carga por fricción en accesorios ocurren en tramos cortos, e hidráulicamente se consideran que ocurren en un punto y usualmente son conocidas como pérdidas de carga localizada o pérdidas menores.

Para cuantificar la pérdida de carga existe un expresión general, que depende de un coeficiente de accesorios,  del diámetro interno y del caudal que circula por el accesorios.

\includegraphics[scale=0.40]{hi.JPG}

A medida que un fluido fluye por un conducto, tubo o algún otro dispositivo, ocurren pérdidas de energía debido a la fricción; tales energías traen como resultado una disminución de la presión entre dos puntos del sistema de flujo. Hay tipos de pérdidas que son muy pequeñas en comparación, y por consiguiente se hace referencia de ellas como pérdidas menores, las cuales ocurren cuando hay un cambio en la sección cruzada de la trayectoria de flujo o en la dirección de flujo, o cuando la trayectoria del flujo se encuentra obstruida como sucede en una válvula, codos, tees, reductores de diámetro, etc
\section{Objetivos}
\subsection{Objetivo General}
\begin{itemize}
\item Determinar experimentalmente los coeficientes de pérdidas por accesorios  para diferentes tipos de accesorios a lo largo de una tubería.
\end{itemize}
\subsection{Objetivo Especifico}
\begin{itemize}
\item Diferenciar los tipos de pérdidas como son por fricción y aditamentos.
\end{itemize}
\begin{itemize}
\item Calcular pérdidas de accesorios la la ecuación de energía de Bernoulli
\end{itemize}
\section{Metodología}
\begin{itemize}
    \item Se identifican las válvulas de control, mediante la manipulación de estas se toman lecturas en el manómetro de mercurio 
\end{itemize}
\begin{itemize}
    \item Se enciende la bomba, la cual da el suministro de agua 
\end{itemize}
\begin{itemize}
    \item Se gradúa la válvula de agua para variar el caudal 
\end{itemize}
\begin{itemize}
    \item Se procede a abrir la válvula correspondiente a cada piezómetro
\end{itemize}
\begin{itemize}
    \item Cada vez que se cambie el accesorio, cierre la válvula que regula el flujo  y apagar la bomba.
\end{itemize}
\begin{itemize}
    \item Al terminar cierre la válvula que controla el fluido y apagar circuito eléctrico
\end{itemize}
\section{Procedimiento}
\begin{itemize}
\item Poner en funcionamiento el banco hidráulico 
\end{itemize}
\begin{center}
\includegraphics[scale=0.15]{Bancopruebas.jpeg}
\end{center}
\begin{itemize}
\item Abrir la válvula para llenar el cilindro
\end{itemize}
\begin{itemize}
\item Regular la presión del agua para medir la altura del mercurio en el cilindro (la altura esta en mm Hg)
\end{itemize}
\begin{center}
\includegraphics[scale=0.15]{Manometro.jpeg}
\end{center}
\begin{itemize}
\item Anotar datos iniciales de la columna de mercurio
\end{itemize}
\begin{itemize}
\item Cerrar la válvulas de los tramos que no son de interés y ubicar el accesorio que se va a usar
\end{itemize}
\begin{itemize}
\item En la primera práctica usar codo de 45 
\end{itemize}
\begin{itemize}
\item Anotar para cada sección del accesorio los valores de presión con los 2 manometros variando el caudal 5 veces tomando nota del tiempo con el cronómetro.
\end{itemize}
\begin{itemize}
\item Iniciar con flujos altos e ir disminuyendo el caudal, tomando los datos de presiones, volúmenes y tiempos.
\end{itemize}
\begin{itemize}
\item Realizar el procedimiento con la unión en Y por lado recto y codo a 90 por radio largo.
\end{itemize}
\section{Discusión de Resultados}
\begin{center}
\includegraphics[scale=0.45]{codo45.JPG}
\end{center}
\begin{center}
\includegraphics[scale=0.45]{uniony.JPG}
\end{center}
\begin{center}
\includegraphics[scale=0.45]{radiolargo.JPG}
\end{center}
\subsection{Velocidad}
La velocidad se obtuvo según los datos del caudal y el diametro de la tuberia
\begin{center}
\includegraphics[scale=0.48]{velocidadcodo45.JPG}
\end{center}
\begin{center}
\includegraphics[scale=0.48]{Velocidaduniony.JPG}
\end{center}
\begin{center}
\includegraphics[scale=0.48]{Velocidadradiolargo.JPG}
\end{center}
\subsection{Perdida de carga H}
Se obtuvo la perdida de carga en mmHg, la diferencia entre las dos alturas se convierte a unidades de mca.
\begin{center}
\includegraphics[scale=0.45]{perdidacodo45.JPG}
\end{center}
\begin{center}
\includegraphics[scale=0.45]{perdidauniony.JPG}
\end{center}
\begin{center}
\includegraphics[scale=0.45]{perdidaradiolargo.JPG}
\end{center}
\subsection{Coeficiente de perdida Ka}
Utilizando la formula de perdida por accesorios, y su respectivo despeje del coeficiente K, se determinan los siguientes valores para cada toma de datos.
\begin{center}
\includegraphics[scale=0.65]{Kacodo45.JPG}
\end{center}
\begin{center}
\includegraphics[scale=0.65]{Kauniony.JPG}
\end{center}
\begin{center}
\includegraphics[scale=0.65]{Karadiolargo.JPG}
\end{center}
\subsection{Promedio Estadictico Ka}
\begin{center}
\includegraphics[scale=0.8]{PromedioKa.JPG}
\end{center}
\subsection{Graficas H/V2/2g}
\begin{center}
    \includegraphics[scale=0.5]{graficacodo45.JPG}
\end{center}
\begin{center}
    \includegraphics[scale=0.5]{graficauniony.JPG}
\end{center}
\begin{center}
    \includegraphics[scale=0.5]{graficaradiolargo.JPG}
\end{center}
\section{Conclusión}
\begin{itemize}
\item La práctica nos da a conocer factores con los que el fluido puede ser transportado y calculado de distintas maneras por la velocidad y presión que dependen de los accesorios.
\end{itemize}
\begin{itemize}
\item Las pérdidas por accesorios son directamente proporcional a la carga de velocidad.
\end{itemize}
\begin{itemize}
\item Las pérdidas por fricción son directamente proporcional a la velocidad del fluido ya que entre mas velocidad mayor son las pérdidas de energía.
\end{itemize}
\section{Bibliografía}
\begin{itemize}
\item \url{https://es.slideshare.net/millos1791/aditamentos-entrega}
\end{itemize}
\begin{itemize}
\item \url{http://mcjabe.blogspot.com/2013/12/informe-perdidas-por-friccion-y.html}
\end{itemize}
\begin{itemize}
\item \url{http://laboratoriosmecanicadefluidos.blogspot.com}
\end{itemize}
\end{multicols}
\end{document}